\documentclass[11pt]{article}
\usepackage[T1]{fontenc}
\usepackage[utf8]{inputenc}
\usepackage[a4paper,margin=1in]{geometry}
\usepackage{titlesec}
\usepackage{enumitem}
\setlist[itemize]{itemsep=0pt,parsep=0pt,topsep=0pt,partopsep=0pt}

% Compact outline headings.
\titleformat{\section}[hang]{\normalfont\normalsize\bfseries}{\thesection}{0.8em}{}
\titleformat{\subsection}[hang]{\normalfont\normalsize}{\thesubsection}{0.8em}{}
\titlespacing*{\section}{0pt}{0pt}{0.3ex}
% Explicit spacing also applies after a subsection with no body text.
\AddToHook{cmd/section/before}{\par\addvspace{3ex plus .5ex minus .3ex}}
\titlespacing*{\subsection}{1.5em}{0pt}{0pt}
\title{Mathematics -- Arithmetic}
\author{}
\date{}

\begin{document}
\maketitle


\section{Numbers and Basic Operations}
\subsection{Numbers: The Building Blocks of Mathematics}
\subsection{The Number Line: Putting Numbers in Their Place}
\subsection{Negative Numbers}
\subsection{Number Sets: Basics}
\subsection{Basic Operations}
\begin{itemize}
    \item Addition
    \item Subtraction
    \item Addition and Subtraction: Signed Integers
    \item Multiplication
    \item Division
    \item Multiplication and Division: Signed Integers
    \item Additive and Multiplicative Identities
\end{itemize}
\subsection{Order of Operations}
\subsection{Absolute Value}

\section{Basic Numeracy and Number Theory}
\subsection{Divisibility}
\subsection{Factors}
\subsection{Multiples}
\subsection{Prime Numbers and Composite Numbers}
\subsection{Product of Primes and Prime Factorization}
\subsection{Least/Lowest Common Multiple (LCM)}
\subsection{Highest/Greatest Common Factor/Divisor (HCF, GCF, HCD)}

\section{Fractions: Parts of a Whole}
\subsection{Introduction to Fractions}
\subsection{Reducing Fractions}
\subsection{Equivalent Fractions}
\subsection{Division by Zero}
\begin{itemize}
    \item Trivia/Extra: Division by Zero
\end{itemize}
\subsection{Arithmetic on Fractions}
\begin{itemize}
    \item Addition and Subtraction
    \item Multiplication and Division
\end{itemize}
\subsection{Signs of Fractions}
\subsection{Reciprocals}

\section{Mixed Numbers and Improper Fractions: Wholes and Parts}
\subsection{Definitions}
\begin{itemize}
    \item Mixed Numbers
    \item Improper Fractions
\end{itemize}
\subsection{Arithmetic on Mixed Numbers and Improper Fractions}
\begin{itemize}
    \item Addition and Subtraction
    \item Multiplication and Division
\end{itemize}
\subsection{Conversion Between Mixed Numbers and Improper Fractions}
\subsection{Relationships between Fractions: Size Comparison and Parts of a Whole}
\subsection{Mixed Measures}

\section{Decimals}
\subsection{Place Values}
\subsection{Arithmetic with Decimals}
\subsection{Decimals and Fractions: Two Sides of the Same Coin}
\subsection{Repeating and Terminating Decimals}
\subsection{Rounding}
\subsection{Percentages}
\begin{itemize}
    \item Percentages as Fractions and Decimals
    \item Conversion Between Percentages, Fractions, and Decimals
    \item Percent of a Number
    \item Percentage Increase and Decrease (Practical Calculations)
\end{itemize}

\section{Ratios and Proportions}
\subsection{Unit Price}
\subsection{Unit Multipliers / Conversion Factors}

\section{Exponents}
\subsection{Rules of Exponents}
\begin{itemize}
    \item Addition and Subtraction
    \item Multiplication and Division: Product and Quotient Rules of Exponents
    \item Power Rule of Exponents
\end{itemize}
\subsection{Negative and Zero Exponents}

\section{Radicals}
\subsection{Rules of Radicals}
\begin{itemize}
    \item Addition/Subtraction
    \item Multiplication and Division: Product and Quotient Rules of Radicals
\end{itemize}
\subsection{Radical Expressions}

\section{Scientific Notation}
\subsection{Powers of 10}
\subsection{Scientific Notation}
\subsection{Multiplication/Division}
\subsection{Estimation}

\section{Basic Measurement}
\subsection{Metric System and Imperial System}
\subsection{Length: Imperial vs Metric}
\subsection{Mass: Imperial vs Metric}
\subsection{Temperature: The Three Scales}
\subsection{Time: The Universal Constant}

\end{document}
